% Preamble template for Tufte-style applied econometrics lecture notes.
% Do not compile this file directly; compile the main notes file instead.
\documentclass{tufte-handout}

\usepackage[T1]{fontenc}
\usepackage[utf8]{inputenc}
\usepackage{ntheorem}
\usepackage{graphicx}
\usepackage{amsmath}
\usepackage{amssymb}
\usepackage{mathtools}
\usepackage{hyperref}
\usepackage{epigraph}
\usepackage{booktabs}
\usepackage{array}
\theoremstyle{break}

\hypersetup{
  colorlinks,
  urlcolor = blue,
  pdfauthor={Swapnil Singh},
  pdfkeywords={econometrics, causal inference, applied econometrics},
  pdftitle={Lecture Notes for Applied Econometrics},
  pdfpagemode=UseNone
}

\newtheorem{ruleN}{Rule}
\newtheorem{thmN}{Theorem}
\newtheorem{assN}{Assumption}
\newtheorem{defN}{Definition}
\newtheorem{exmp}{Example}
\newtheorem{cmt}{Comment}
\newtheorem{proof}{Proof}

\newcommand{\continuation}{??}
\newtheorem*{excont}{Example \continuation}
\newenvironment{continueexample}[1]
 {\renewcommand{\continuation}{\ref{#1}}\excont[continued]}
 {\endexcont}

\newcommand{\bY}{\mathbf{Y}}
\newcommand{\bX}{\mathbf{X}}
\newcommand{\bD}{\mathbf{D}}
\newcommand{\E}{\mathbb{E}}
\newcommand{\Prob}{\mathbb{P}}
\newcommand{\Var}{\mathrm{Var}}
\newcommand{\Cov}{\mathrm{Cov}}
\newcommand{\se}{\mathrm{SE}}
\newcommand{\CI}{\mathrm{CI}}
\newcommand{\ATE}{\mathrm{ATE}}
\newcommand{\ATT}{\mathrm{ATT}}
\newcommand{\TOT}{\mathrm{TOT}}
\newcommand{\LATE}{\mathrm{LATE}}
\newcommand{\ITT}{\mathrm{ITT}}
\newcommand{\RD}{\mathrm{RD}}
\newcommand{\DiD}{\mathrm{DiD}}
\newcommand{\MMSE}{\mathrm{MMSE}}
\newcommand{\OLS}{\mathrm{OLS}}
\newcommand{\Normal}{\mathcal{N}}
\newcommand{\plim}{\operatorname*{plim}}
\DeclareMathOperator*{\argmin}{arg\,min}
\DeclareMathOperator{\Supp}{Supp}

\newcommand\independent{\protect\mathpalette{\protect\independenT}{\perp}}
\def\independenT#1#2{\mathrel{\rlap{$#1#2$}\mkern2mu{#1#2}}}
\newcommand{\indep}{\perp\!\!\!\perp}

\usepackage{cleveref}
\crefname{appsec}{appendix}{appendices}
\crefname{appsubsec}{appendix}{appendices}
\crefname{assumption}{assumption}{assumptions}
\crefname{equation}{equation}{equations}
\crefname{exmp}{example}{examples}
\crefname{assN}{assumption}{assumptions}
\crefname{cmt}{comment}{comments}
\crefname{defN}{definition}{definitions}

\usepackage[nolist]{acronym}
\begin{acronym}
  \acro{ATE}{average treatment effect}
  \acro{ATT}{average treatment effect on the treated}
  \acro{CIA}{conditional independence assumption}
  \acro{CI}{confidence interval}
  \acro{CLT}{central limit theorem}
  \acro{DiD}{difference-in-differences}
  \acro{FE}{fixed effects}
  \acro{FWL}{Frisch-Waugh-Lovell}
  \acro{IV}{instrumental variables}
  \acro{LATE}{local average treatment effect}
  \acro{OLS}{ordinary least squares}
  \acro{OVB}{omitted-variable bias}
  \acro{PO}{potential outcomes}
  \acro{RCT}{randomized controlled trial}
  \acro{RD}{regression discontinuity}
  \acro{SUTVA}{stable unit treatment value assignment}
  \acro{TWFE}{two-way fixed effects}
\end{acronym}

\def\inprobHIGH{\,{\buildrel p \over \rightarrow}\,}
\def\inprob{\,{\inprobHIGH}\,}
\def\indistHIGH{\,{\buildrel d \over \rightarrow}\,}
\def\indist{\,{\indistHIGH}\,}

\usepackage[many]{tcolorbox}
\definecolor{main}{HTML}{3F6EA7}
\definecolor{sub}{HTML}{F2F7FD}
\definecolor{sub2}{HTML}{FFF5E8}
\definecolor{mainline}{HTML}{D3DFEE}
\definecolor{warn}{HTML}{C7771B}
\definecolor{warnline}{HTML}{EAD9C0}

\tcbset{
    enhanced,
    breakable,
    colback = white,
    boxrule = 0.5pt,
    arc = 2pt,
    left = 9pt,
    right = 9pt,
    top = 7pt,
    bottom = 7pt,
    before skip = 0.3cm,
    after skip = 0.45cm
}

\newtcolorbox{boxD}{
    colback = sub,
    colframe = mainline,
    borderline west = {2.2pt}{0pt}{main},
    borderline north = {0.6pt}{0pt}{main!25},
    borderline south = {0.6pt}{0pt}{main!15}
}

\newtcolorbox{boxF}{
    colback = sub2,
    colframe = warnline,
    borderline west = {2.2pt}{0pt}{warn},
    borderline north = {0.6pt}{0pt}{warn!25},
    borderline south = {0.6pt}{0pt}{warn!15}
}

\usepackage{colortbl}
\usepackage{tikz}
\usepackage{verbatim}
\usetikzlibrary{positioning}
\usetikzlibrary{snakes}
\usetikzlibrary{calc}
\usetikzlibrary{arrows}
\usetikzlibrary{decorations.markings}
\usetikzlibrary{shapes.misc}
\usetikzlibrary{matrix,shapes,arrows,fit,tikzmark}


\title{Applied Econometrics: Session 1 Notes}
\author{Swapnil Singh}
\date{}

\hypersetup{
  pdftitle={Applied Econometrics: Session 1 Notes},
  pdfauthor={Swapnil Singh},
  pdfkeywords={causal inference, potential outcomes, experiments, selection bias, identification}
}

\setlength{\headheight}{26pt}

\begin{document}
\maketitle

\begin{abstract}
This lecture builds the language used in the rest of the course. We start from
causal questions rather than estimators. A causal question is a question about a
counterfactual: what would have happened to the same unit under a different
treatment status? We then introduce potential outcomes, the fundamental problem
of causal inference, selection bias, random assignment, treatment effects from
means and regressions, and the discipline of using an ideal experiment to guide
empirical work.
\end{abstract}

\begin{boxD}
\textbf{Reading and objective.} These notes accompany Session 1, based on the
first part of the course material from Mostly Harmless Econometrics and the
experimental ideal in modern causal inference. The objective is to learn how to
state a causal question, translate it into potential outcomes, identify the
missing counterfactual, and explain why random assignment is the benchmark
against which all later quasi-experimental designs are judged.
\end{boxD}

\section{Start with the Question}

Applied econometrics does not begin with a regression command. It begins with a
question. The command is useful only after we know what object the command is
supposed to estimate. Many weak empirical projects fail because the estimator is
clear but the causal question is vague.

\begin{boxF}
\textbf{Common beginner mistake.} A student says, ``I want to regress wages on
schooling.'' That is not yet a research question. It is a proposed calculation.
The research question is something like: ``What is the causal effect of one more
year of schooling on wages for a specified population?''
\end{boxF}

\subsection{The Four Research FAQs}

For any applied project, ask four questions in order.
\begin{enumerate}
    \item What is the causal relationship of interest?
    \item What ideal experiment would answer it?
    \item What real-world variation identifies the effect in the available data?
    \item What mode of statistical inference matches the sampling or assignment
    process?
\end{enumerate}

These are sequential. If the first question is unclear, the rest of the project
will be confused. If the ideal experiment is not stated, it is hard to say what
the observational design is trying to mimic. If the identification strategy is
not credible, standard errors cannot rescue the estimate. If inference ignores
the design, the estimate may look much more precise than it really is.

\begin{defN}[Causal relationship]
A causal relationship describes how an outcome for the same unit would change
if treatment were changed, holding fixed the relevant background conditions.
It is not merely an association between treated and untreated units.
\end{defN}

Consider schooling and wages. The descriptive fact is that workers with more
schooling often earn more. The causal question is different: what would happen
to a worker's wage if that same worker received more schooling than she
otherwise would have received? The difference matters because workers with more
schooling may differ in ability, family background, motivation, location, or
health even before schooling decisions are made.

\subsection{The Ideal Experiment}

The ideal experiment is a thought experiment. We imagine a randomized trial with
perfect compliance, no attrition, no spillovers, and clean measurement. The
point is not that the experiment is feasible. The point is that the ideal
experiment forces us to define treatment, outcome, population, timing, and the
counterfactual.

\begin{exmp}[Returns to schooling]
Suppose the target is the effect of one extra year of schooling on adult wages.
The ideal experiment would randomly assign otherwise comparable students to
receive one additional year of schooling, prevent all other channels from
changing, and compare adult wages. In practice we cannot force schooling in this
way. But the thought experiment teaches us what later designs, such as
compulsory schooling laws, are trying to approximate.
\end{exmp}

\begin{boxD}
\textbf{Design-first habit.} Before naming an estimator, write a sentence of
the form: ``The ideal experiment would randomly assign ... and compare ... .''
If that sentence cannot be written, the empirical design is not yet ready.
\end{boxD}

\subsection{Fundamentally Unidentified Questions}

Some questions are hard not because data are missing, but because the causal
object itself cannot be separated cleanly. The slides call these
\emph{fundamentally unidentified questions}.

\begin{exmp}[School starting age]
Suppose we ask for the effect of being older at school entry on first-grade test
scores. A child who starts school later is older when observed and may also have
different time in school depending on the comparison date. Age, grade, and
school exposure are mechanically linked. Comparing children at the same test
date changes age; comparing at the same age changes schooling exposure. Without
a sharper question, there is no single separable treatment effect.
\end{exmp}

The lesson is not that all difficult questions should be abandoned. The lesson
is that we must be precise about what is being varied. If the treatment bundle
contains age, schooling exposure, grade, peer group, and parental choices, then
the estimated effect is an effect of that bundle. Calling it the effect of
``age'' alone may be misleading.

\section{Potential Outcomes}

Potential outcomes give formal notation to the counterfactual idea. For each
unit \(i\), define
\[
    Y_{1i} = \text{outcome if treated}, \qquad
    Y_{0i} = \text{outcome if untreated}.
\]
Let \(D_i\in\{0,1\}\) be the treatment indicator. The observed outcome is
\[
    Y_i = D_iY_{1i} + (1-D_i)Y_{0i}
        = Y_{0i} + D_i(Y_{1i}-Y_{0i}).
\]

\begin{defN}[Individual causal effect]
The individual causal effect for unit \(i\) is
\[
    Y_{1i}-Y_{0i}.
\]
It compares two potential outcomes for the same unit.
\end{defN}

This object is conceptually simple but empirically impossible to observe. If a
person receives the treatment, we observe \(Y_{1i}\) and miss \(Y_{0i}\). If the
person does not receive the treatment, we observe \(Y_{0i}\) and miss \(Y_{1i}\).

\begin{boxF}
\textbf{Fundamental problem.} For each unit at a point in time, we observe only
one potential outcome. Causal inference is therefore a missing-counterfactual
problem, not merely a data-size problem.
\end{boxF}

\subsection{Average Treatment Effects}

Since individual treatment effects are not fully observed, applied work usually
targets averages.
\[
    \ATE = \E[Y_{1i}-Y_{0i}]
\]
is the average treatment effect in the full population. The average treatment
effect on the treated is
\[
    \ATT = \E[Y_{1i}-Y_{0i}\mid D_i=1].
\]

\begin{boxD}
\textbf{ATE versus ATT.} The ATE asks what would happen if treatment status
changed for everyone in the target population. The ATT asks what the treatment
did for the units who actually received it. These are the same only under
homogeneous treatment effects or special sampling and assignment conditions.
\end{boxD}

The distinction matters. A job-training program may help disadvantaged workers
who enroll but not help all workers equally. A medical treatment may have a
large effect for patients who choose it because they are ill, but a smaller
average effect in the entire population.

\section{Selection Bias}

The observed difference in mean outcomes between treated and untreated units is
\[
    \E[Y_i\mid D_i=1]-\E[Y_i\mid D_i=0].
\]
This is easy to compute, but it is not automatically causal. Using the potential
outcomes equation,
\[
    \E[Y_i\mid D_i=1]=\E[Y_{1i}\mid D_i=1],
\]
and
\[
    \E[Y_i\mid D_i=0]=\E[Y_{0i}\mid D_i=0].
\]
Therefore
\[
\begin{aligned}
    \E[Y_i\mid D_i=1]-\E[Y_i\mid D_i=0]
    &=
    \E[Y_{1i}\mid D_i=1]-\E[Y_{0i}\mid D_i=0] \\
    &=
    \underbrace{\E[Y_{1i}-Y_{0i}\mid D_i=1]}_{\ATT}
    +
    \underbrace{\E[Y_{0i}\mid D_i=1]-\E[Y_{0i}\mid D_i=0]}_{\text{selection bias}}.
\end{aligned}
\]

\begin{thmN}[Selection-bias decomposition]
The observed treated-minus-untreated difference equals the treatment effect on
the treated plus the baseline untreated-outcome difference between treated and
untreated units.
\end{thmN}

\begin{proof}
Add and subtract \(\E[Y_{0i}\mid D_i=1]\):
\[
\begin{aligned}
    \E[Y_{1i}\mid D_i=1]-\E[Y_{0i}\mid D_i=0]
    &=
    \E[Y_{1i}\mid D_i=1]-\E[Y_{0i}\mid D_i=1] \\
    &\quad+
    \E[Y_{0i}\mid D_i=1]-\E[Y_{0i}\mid D_i=0] \\
    &=
    \E[Y_{1i}-Y_{0i}\mid D_i=1]
    +
    \E[Y_{0i}\mid D_i=1]-\E[Y_{0i}\mid D_i=0].
\end{aligned}
\]
\end{proof}

\begin{exmp}[Hospitals and health]
Suppose lower values of \(Y\) mean better health. Hospital patients have an
average outcome of \(4.3\), while nonpatients have an average outcome of \(3.6\).
The naive difference is \(0.7\), suggesting hospital care is associated with
worse health. But hospital patients were probably sicker before treatment. If
their untreated potential outcome would have been \(1.0\) point worse than that
of nonpatients, then the causal effect for the treated is
\[
    0.7 - 1.0 = -0.3.
\]
The naive comparison points in the wrong direction.
\end{exmp}

\section{Random Assignment}

Random assignment is the benchmark because it removes selection bias by design.
If treatment is randomly assigned, then treatment status is independent of the
potential outcomes:
\[
    (Y_{1i},Y_{0i}) \indep D_i.
\]

\begin{assN}[Random assignment]
Treatment is randomly assigned if
\[
    (Y_{1i},Y_{0i}) \indep D_i.
\]
Equivalently, treated and untreated groups have the same distribution of
potential outcomes before treatment is realized.
\end{assN}

Under random assignment,
\[
    \E[Y_{0i}\mid D_i=1]=\E[Y_{0i}\mid D_i=0],
\]
so the selection-bias term is zero. Also,
\[
    \E[Y_{1i}\mid D_i=1]=\E[Y_{1i}\mid D_i=0]=\E[Y_{1i}].
\]
Therefore the observed mean difference identifies the ATE:
\[
    \E[Y_i\mid D_i=1]-\E[Y_i\mid D_i=0]
    =
    \E[Y_{1i}]-\E[Y_{0i}]
    =
    \ATE.
\]

\begin{thmN}[Identification by random assignment]
If treatment is randomly assigned and outcomes are observed for all assigned
units, then the treated-minus-control mean difference identifies the average
treatment effect.
\end{thmN}

\begin{boxF}
\textbf{What experiments do not automatically solve.} Random assignment solves
selection into treatment at baseline. It does not automatically solve
non-compliance, attrition, spillovers, bad measurement, small sample noise, or
incorrect standard errors.
\end{boxF}

\subsection{The Tennessee STAR Example}

The Tennessee STAR class-size experiment is useful because it shows what good
experimental evidence looks like in practice. Students and teachers were
assigned to classes of different sizes. The core empirical question is whether
smaller classes improved achievement.

The first diagnostic is balance. Before looking at outcomes, check whether
pre-treatment characteristics are similar across treatment groups. Balance is
not proof that randomization occurred, but it is evidence that the randomization
created comparable groups in the observed data.

\begin{boxD}
\textbf{Balance interpretation.} In a randomized experiment, small baseline
differences can happen by chance. The question is whether the differences are
large, systematic, or suspicious relative to the assignment process.
\end{boxD}

\section{Experiments as Regressions}

With a binary treatment and an intercept, the regression
\[
    Y_i = \alpha + \tau D_i + u_i
\]
is just a difference in means. The fitted mean for controls is \(\alpha\). The
fitted mean for treated units is \(\alpha+\tau\). Therefore
\[
    \tau
    =
    \E[Y_i\mid D_i=1]-\E[Y_i\mid D_i=0]
\]
in the population, and the sample OLS estimate equals the sample treated-control
difference.

\begin{thmN}[Binary-treatment OLS]
In a regression of \(Y_i\) on an intercept and a binary treatment \(D_i\), the
OLS slope on \(D_i\) is the treated sample mean minus the control sample mean.
\end{thmN}

\begin{proof}
OLS residuals sum to zero within each value of the binary regressor. Hence the
fitted value for \(D_i=0\) equals the sample mean among controls and the fitted
value for \(D_i=1\) equals the sample mean among treated units. Since these
fitted values are \(\hat{\alpha}\) and \(\hat{\alpha}+\hat{\tau}\), we have
\[
    \hat{\tau}=\bar{Y}_1-\bar{Y}_0.
\]
\end{proof}

Adding pre-treatment covariates to an experimental regression can improve
precision and adjust for chance imbalance:
\[
    Y_i=\alpha+\tau D_i+X_i'\gamma+u_i.
\]
In a correctly randomized experiment, the causal credibility comes from
assignment, not from controls. Controls are useful mainly because they explain
outcome variation and can make the estimate more precise.

\section{From Experiments to Quasi-Experiments}

Most empirical work is not a clean randomized trial. Later lectures study ways
to approximate an experiment using observational data: regression adjustment,
instrumental variables, matching, difference-in-differences, synthetic control,
and regression discontinuity.

The common logic is always the same:
\begin{enumerate}
    \item define the missing counterfactual;
    \item state what comparison group is being used to approximate it;
    \item defend why that comparison group is credible;
    \item report uncertainty honestly.
\end{enumerate}

\begin{boxD}
\textbf{Course map.} Every later method is a way of answering the same
question: why should the observed comparison group reveal the missing potential
outcome?
\end{boxD}

\section{Worked Problems}

\begin{exmp}[Observed difference and selection bias]
An employment program has average earnings of 24 among participants and 20 among
nonparticipants. Suppose participants would have earned 18 without the program,
while nonparticipants would have earned 20 without the program. The observed
difference is
\[
    24-20=4.
\]
The ATT is
\[
    24-18=6.
\]
The selection-bias term is
\[
    18-20=-2.
\]
Thus
\[
    4=6+(-2).
\]
The naive comparison understates the effect because participants were negatively
selected on untreated earnings.
\end{exmp}

\begin{exmp}[Random assignment]
In an experiment, the treated mean is 75 and the control mean is 70. Under
random assignment, the ATE estimate is
\[
    \hat{\tau}=75-70=5.
\]
The reason this has a causal interpretation is not the arithmetic. The reason is
that random assignment makes the control mean an estimate of the treated group's
missing untreated mean.
\end{exmp}

\begin{exmp}[Regression with a treatment dummy]
Suppose three treated observations have outcomes \(8,10,12\), and three control
observations have outcomes \(5,6,7\). The treated mean is \(10\), the control
mean is \(6\), and the regression slope in
\[
    Y_i=\alpha+\tau D_i+u_i
\]
is
\[
    \hat{\tau}=10-6=4.
\]
The intercept is \(\hat{\alpha}=6\), and the fitted treated mean is
\(\hat{\alpha}+\hat{\tau}=10\).
\end{exmp}

\section{Checklist for Students}

After this lecture you should be able to:
\begin{itemize}
    \item distinguish a causal question from a regression specification;
    \item state the four research FAQs in order;
    \item describe the ideal experiment for a causal question;
    \item define \(Y_{1i}\), \(Y_{0i}\), \(D_i\), \(Y_i\), ATE, and ATT;
    \item explain the fundamental problem of causal inference;
    \item derive the selection-bias decomposition;
    \item explain why random assignment removes selection bias;
    \item compute treatment effects from treatment and control means;
    \item explain why a binary-treatment regression equals a difference in means;
    \item list problems that random assignment does not automatically solve;
    \item connect later quasi-experimental methods to the experimental ideal.
\end{itemize}

\section{Practice Questions}

\begin{enumerate}
    \item Write a causal question about schooling and wages. Then write the
    corresponding ideal experiment.
    \item Explain why ``people who go to hospital are less healthy than people
    who do not'' is not evidence that hospitals harm health.
    \item Define \(Y_{1i}\), \(Y_{0i}\), and \(D_i\) for a job-training program.
    \item What is the individual treatment effect? Why is it not observed?
    \item Derive the selection-bias decomposition by adding and subtracting
    \(\E[Y_{0i}\mid D_i=1]\).
    \item Give an example where selection bias is positive. Give another where
    it is negative.
    \item State the random-assignment condition in potential-outcomes notation.
    \item Explain why random assignment makes the control group useful.
    \item In an experiment, treated mean outcome is 12 and control mean outcome
    is 9. What is the estimated effect?
    \item Show that a regression on an intercept and a treatment dummy gives the
    same number as the difference in means.
    \item Why might covariates still be useful in an experiment?
    \item Name two problems that can remain even in a randomized experiment.
    \item What is a fundamentally unidentified question? Give an example.
    \item Why should an empirical paper describe the design before the
    estimator?
    \item Choose one later method from the course and explain how it tries to
    mimic an experiment.
\end{enumerate}

\end{document}
